% GENERATED FILE. Edit canonical lessons, then run npm run generate:books.
% canonical-source-hash: fnv1a64:870515d4f435923c
% canonical-lessons: PT-C493-bisavo, PT-C493-bisneto, PT-C493-orfao, PT-C493-adotar, PT-C493-adocao

\chapter{Bisavós, bisnetos e adoção}
\label{ch:pt-bisavos-bisnetos-e-adocao}
\hlchaptermodality{493}

\begin{chapteropening}
\textbf{By the end of this chapter:} \emph{I can talk about great-grandparents, great-grandchildren, orphans and adoption.}

Complete the last lesson of chapter 493: I can talk about great-grandparents, great-grandchildren, orphans and adoption.
\end{chapteropening}

\section[o bisavô]{o bisavô — great-grandfather}

\label{lesson:PT-C493-bisavo}



\begin{quote}
Before the new one: say the Portuguese for stepson, then the Portuguese for relative.
\end{quote}

\subsection*{You'll want to know: o bisavô}
\textbf{o bisavô} — \textquotedblleft{}great-grandfather\textquotedblright{}.

\begin{quote}
O meu bisavô nasceu numa aldeia perto do mar.
\end{quote}

\begin{quote}
\emph{My great-grandfather was born in a village near the sea.}
\end{quote}

\emph{A bisavó} — great-grandmother; \emph{os bisavós} — great-grandparents.

\begin{cousinweb}
\emph{bisavô} is built from \emph{bis-} \textquotedblleft{}twice\textquotedblright{} and \emph{avô} \textquotedblleft{}grandfather\textquotedblright{}, which goes back to Latin \emph{avus} \textquotedblleft{}grandfather\textquotedblright{} (Lewis \& Short).
\end{cousinweb}

\begin{grammarlens}[title={What you've built}]
One more word for telling, explaining and describing at B1.
\end{grammarlens}

\subsection*{Your turn}
\begin{itemize}
\raggedright
  \item \emph{Say it:} \emph{o bisavô}
  \item \emph{Say it:} \emph{o bisavô}, once more
  \item \emph{Say it:} a sentence of your own with \emph{o bisavô}
  \item \emph{Recall:} say the Portuguese for stepson, then the Portuguese for relative, then say \emph{o bisavô} again
\end{itemize}

\begin{tcolorbox}[breakable,colback=teal!4,colframe=teal!35!black,title={Before you move on}]
What does o bisavô mean? (\textquotedblleft{}Great-grandfather\textquotedblright{}.) Say it once more.
\end{tcolorbox}
\section[o bisneto]{o bisneto — great-grandson}

\label{lesson:PT-C493-bisneto}



\begin{quote}
Before the new one: say the Portuguese for relative, then the Portuguese for great-grandfather.
\end{quote}

\subsection*{You'll want to know: o bisneto}
\textbf{o bisneto} — \textquotedblleft{}great-grandson\textquotedblright{}.

\begin{quote}
A senhora de noventa anos já tem dois bisnetos.
\end{quote}

\begin{quote}
\emph{The ninety-year-old lady already has two great-grandchildren.}
\end{quote}

\emph{A bisneta} — great-granddaughter.

\begin{cousinweb}
\emph{bisneto} is built from \emph{bis-} \textquotedblleft{}twice\textquotedblright{} and \emph{neto} \textquotedblleft{}grandson\textquotedblright{}, which goes back to Latin \emph{nepōs} \textquotedblleft{}grandson\textquotedblright{} (Lewis \& Short).
\end{cousinweb}

\begin{grammarlens}[title={What you've built}]
One more word for telling, explaining and describing at B1.
\end{grammarlens}

\subsection*{Your turn}
\begin{itemize}
\raggedright
  \item \emph{Say it:} \emph{o bisneto}
  \item \emph{Say it:} \emph{o bisneto}, once more
  \item \emph{Say it:} a sentence of your own with \emph{o bisneto}
  \item \emph{Recall:} say the Portuguese for relative, then the Portuguese for great-grandfather, then say \emph{o bisneto} again
\end{itemize}

\begin{tcolorbox}[breakable,colback=teal!4,colframe=teal!35!black,title={Before you move on}]
What does o bisneto mean? (\textquotedblleft{}Great-grandson\textquotedblright{}.) Say it once more.
\end{tcolorbox}
\section[o órfão]{o órfão — orphan}

\label{lesson:PT-C493-orfao}



\begin{quote}
Before the new one: say the Portuguese for great-grandfather, then the Portuguese for great-grandson.
\end{quote}

\subsection*{You'll want to know: o órfão}
\textbf{o órfão} — \textquotedblleft{}orphan\textquotedblright{}.

\begin{quote}
O rapaz ficou órfão muito novo e foi viver com os avós.
\end{quote}

\begin{quote}
\emph{The boy was orphaned very young and went to live with his grandparents.}
\end{quote}

Also an adjective: \emph{ficar órfão} — to lose your parents. A girl is \emph{a órfã}; the plural is \emph{órfãos}.

\begin{cousinweb}
\emph{órfão} is inherited from Late Latin \emph{orphanus} \textquotedblleft{}an orphan\textquotedblright{} (Lewis \& Short), which was taken from Greek \emph{orphanós} \textquotedblleft{}without parents\textquotedblright{} (EtymDB).
\end{cousinweb}

\begin{grammarlens}[title={What you've built}]
One more word for telling, explaining and describing at B1.
\end{grammarlens}

\subsection*{Your turn}
\begin{itemize}
\raggedright
  \item \emph{Say it:} \emph{o órfão}
  \item \emph{Say it:} \emph{o órfão}, once more
  \item \emph{Say it:} a sentence of your own with \emph{o órfão}
  \item \emph{Recall:} say the Portuguese for great-grandfather, then the Portuguese for great-grandson, then say \emph{o órfão} again
\end{itemize}

\begin{tcolorbox}[breakable,colback=teal!4,colframe=teal!35!black,title={Before you move on}]
What does o órfão mean? (\textquotedblleft{}Orphan\textquotedblright{}.) Say it once more.
\end{tcolorbox}
\section[adotar]{adotar — to adopt}

\label{lesson:PT-C493-adotar}



\begin{quote}
Before the new one: say the Portuguese for great-grandson, then the Portuguese for orphan.
\end{quote}

\subsection*{You'll want to know: adotar}
\textbf{adotar} — \textquotedblleft{}to adopt\textquotedblright{}.

\begin{quote}
Eles adotaram uma menina de três anos.
\end{quote}

\begin{quote}
\emph{They adopted a three-year-old girl.}
\end{quote}

Before AO90 spelled \emph{adoptar} in Portugal.

\begin{cousinweb}
\emph{adotar} goes back, through Old Portuguese \emph{adoutar}, to Latin \emph{adoptāre} \textquotedblleft{}to choose for oneself, to adopt\textquotedblright{} (EtymDB), built on \emph{optāre} \textquotedblleft{}to choose\textquotedblright{}.
\end{cousinweb}

\begin{grammarlens}[title={What you've built}]
One more word for telling, explaining and describing at B1.
\end{grammarlens}

\subsection*{Your turn}
\begin{itemize}
\raggedright
  \item \emph{Say it:} \emph{adotar}
  \item \emph{Say it:} \emph{adotar}, once more
  \item \emph{Say it:} a sentence of your own with \emph{adotar}
  \item \emph{Recall:} say the Portuguese for great-grandson, then the Portuguese for orphan, then say \emph{adotar} again
\end{itemize}

\begin{tcolorbox}[breakable,colback=teal!4,colframe=teal!35!black,title={Before you move on}]
What does adotar mean? (\textquotedblleft{}To adopt\textquotedblright{}.) Say it once more.
\end{tcolorbox}
\section[a adoção]{a adoção — adoption}

\label{lesson:PT-C493-adocao}



\begin{quote}
Before the new one: say the Portuguese for orphan, then the Portuguese for to adopt.
\end{quote}

\subsection*{You'll want to know: a adoção}
\textbf{a adoção} — \textquotedblleft{}adoption\textquotedblright{}.

\begin{quote}
A adoção demorou quase dois anos.
\end{quote}

\begin{quote}
\emph{The adoption took almost two years.}
\end{quote}

Before AO90 spelled \emph{adopção} in Portugal.

\begin{cousinweb}
\emph{adoção} comes from Latin \emph{adoptiō} \textquotedblleft{}adoption\textquotedblright{} (EtymDB), from \emph{adoptāre}.
\end{cousinweb}

\begin{grammarlens}[title={What you've built}]
One more word for telling, explaining and describing at B1.
\end{grammarlens}

\subsection*{Your turn}
\begin{itemize}
\raggedright
  \item \emph{Say it:} \emph{a adoção}
  \item \emph{Say it:} \emph{a adoção}, once more
  \item \emph{Say it:} a sentence of your own with \emph{a adoção}
  \item \emph{Recall:} say the Portuguese for orphan, then the Portuguese for to adopt, then say \emph{a adoção} again
\end{itemize}

\begin{tcolorbox}[breakable,colback=teal!4,colframe=teal!35!black,title={Before you move on}]
What does a adoção mean? (\textquotedblleft{}Adoption\textquotedblright{}.) Say it once more.
\end{tcolorbox}
